\section{Conceptualization}

In this section we present an analysis of the process that led to expressing the
knowledge in logical form. In particular, we discuss the objects that were
identified and the relations and properties of those objects.

The knowledge is organised over four inference levels, which structure the
reasoning from the most abstract diagnosis down to the elementary observations:

\begin{enumerate}[label=(\Roman*)]
    \item \emph{Disorders} ---the ``child'' anxiety disorders that the system is
    able to diagnose, encoded by the relation \texttt{disorder/3};
    \item \emph{Sub-disorders and types} ---the ``parent'' or
    ``non-codifiable'' disorders and the specialisation types that, combined with
    them, define a disorder, encoded by the relation \texttt{disorder\_nc/4};
    \item \emph{Symptomatic manifestations and features} ---the complex symptoms
    and the features attached to them, encoded by the relations
    \texttt{symptomatic\_manifestation/4} and \texttt{symptomatic\_feature/4};
    \item \emph{Elementary symptoms} ---the atomic facts reported by the user,
    encoded by the relation \texttt{symptom/3}.
\end{enumerate}


\subsection{Objects}

The conceptualization phase brought several objects to light. These objects can be
classified into categories.

In Table~\ref{tab:objects}, rather than weighing down the discussion with a list
of every single identified object, we prefer to report the categories to which the
objects belong, focusing attention on the description of what they represent. For
each category, a few examples of objects are also given.

\begin{longtable}{p{0.22\linewidth} p{0.42\linewidth} p{0.28\linewidth}}
\caption{Objects identified during the conceptualization
phase.}\label{tab:objects}\\
\toprule
\textbf{Object Category} & \textbf{Description} & \textbf{Examples} \\
\midrule
\endfirsthead
\toprule
\textbf{Object Category} & \textbf{Description} & \textbf{Examples} \\
\midrule
\endhead
\midrule
\multicolumn{3}{r}{\footnotesize\itshape continued on next page}\\
\endfoot
\bottomrule
\endlastfoot

Disorder &
The ``child'' anxiety disorders that the expert system is able to diagnose. &
Panic Disorder Without Agoraphobia; Specific Phobia, Animal Type; Generalised
Anxiety Disorder; \dots \\
\addlinespace

Non-Codifiable Disorder &
Both the ``parent'' anxiety disorders to which the ``child'' disorders refer and
the so-called ``non-codifiable'' disorders, that is, those disorders that, without
a precise specialisation or placement, cannot be recorded. &
Agoraphobia; Specific Phobia; Post-Traumatic Stress Disorder; Panic Disorder;
Social Phobia; \dots \\
\addlinespace

Type of Non-Codifiable Disorder &
The various specialisations that, combined with the non-codifiable disorders,
allow a disorder to be defined. &
Generalised Type; Circumscribed Type; Animal Type; Natural Environment Type;
\dots \\
\addlinespace

Symptomatic Manifestation &
The various manifestations of symptoms that denote an anxiety disorder.
Typically, a symptomatic manifestation is made up of several criteria (complex
symptoms). Each criterion is composed of one or more symptoms, which need not all
be mandatorily required for the diagnosis: it may be the case that a minimum
number of symptoms out of a larger set is required. &
Increased Arousal; Psychological Trauma; Social Phobic Anxiety; Agoraphobic
Anxiety; Unexpected Panic Attack; Generalised Anxiety; \dots \\
\addlinespace

Symptomatic Feature &
The features connected to the symptoms that allow one disorder to be determined or
distinguished from another. &
Duration of at least 6 months; Frequent panic attacks; Onset of symptoms within 1
month of the traumatic event; \dots \\
\addlinespace

Elementary Symptom &
The atomic symptoms that, once assembled, give rise to the symptomatic
manifestations. &
Depersonalisation; Terror; Avoidance; Nausea; \dots \\
\addlinespace

Patient &
The categories of patient to whom the diagnosis is addressed. Placing the patient
in one of the three age ranges makes it possible to distinguish both the number of
symptomatic manifestations or features required by a given disorder and the number
of symptoms expected from a single symptomatic manifestation. &
Child (0--9 years); Adolescent (10--17 years); Adult (18+). \\
\addlinespace

Degree of Insight of the Obsessions &
The degrees of insight into the reasonableness of the obsessions-compulsions. &
High Insight; Low Insight; \dots \\
\addlinespace

Type of Anxiety &
The types of anxiety. The distinction is based on the main reason that triggers
the anxious process. &
Anxiety due to embarrassment in public; Anxiety due to concern for one's own
health; \dots \\
\addlinespace

Type of Panic Attack &
The two types into which a panic attack can be specialised. &
Unexpected Panic Attack; Situationally Triggered Panic Attack. \\
\addlinespace

Type of Phobic Reaction &
The types of reaction that an individual may have towards a feared situation or
object. &
Endurance with distress; Avoidance; \dots \\
\addlinespace

Symptom Duration &
The various durations used to verify all the symptoms that require the presence of
the symptoms for a precise length of time. &
Less than 2 days; 2 days -- less than 1 month; 1 month -- less than 3 months;
3 months -- less than 6 months; 6 months or more. \\
\addlinespace

Symptom Onset &
The various periods within which the acute stress disorder and the post-traumatic
stress disorder can arise. &
Less than 1 month; 1 month -- less than 6 months; 6 months or more. \\
\addlinespace

Panic Attack Frequency &
The various frequencies of panic attacks required in panic disorder. &
1; 2--5; 6+. \\
\end{longtable}


\subsection{Relations and Properties}

Table~\ref{tab:relations} reports the relations and the properties of the objects.
The first five relations correspond to the four inference levels described at the
beginning of this section.

\begin{longtable}{p{0.46\linewidth} p{0.46\linewidth}}
\caption{Relations and properties of the objects.}\label{tab:relations}\\
\toprule
\textbf{Relation} & \textbf{Description} \\
\midrule
\endfirsthead
\toprule
\textbf{Relation} & \textbf{Description} \\
\midrule
\endhead
\midrule
\multicolumn{2}{r}{\footnotesize\itshape continued on next page}\\
\endfoot
\bottomrule
\endlastfoot

\texttt{patient(Patient)} &
Indicates the type of patient to whom the diagnosis is addressed. \\
\addlinespace

\texttt{disorder/3} &
States that the given object is a ``child'' disorder and is diagnosable for the
specified patient (level~I). \\
\addlinespace

\texttt{disorder\_nc/4} &
States that the given object is a ``parent'' or ``non-codifiable'' disorder,
together with its specialisation type, and is diagnosable for the specified
patient (level~II). \\
\addlinespace

\texttt{symptomatic\_manifestation/4} &
States that the given object is a symptomatic manifestation and is diagnosable for
the specified patient (level~III). \\
\addlinespace

\texttt{symptomatic\_feature/4} &
States that the given object is a symptomatic feature and is diagnosable for the
specified patient (level~III). \\
\addlinespace

\texttt{symptom/3} &
Encodes a criterion (complex symptom) at the elementary level (level~IV): the
list argument indicates the various elementary symptoms that compose the
criterion, while the numeric argument indicates the number of symptoms mandatorily
required to diagnose the criterion. \\
\addlinespace

\texttt{aetiological\_evidence([Evidence\_List], Evidence\_Number)} &
Indicates an aetiological evidence, where \texttt{Evidence\_List} lists the
various ways of verifying the aetiological link between the disorder and the
factor under examination (drug, withdrawal, intoxication), and
\texttt{Evidence\_Number} indicates the number of evidences mandatorily required
for the diagnosis. \\
\addlinespace

\texttt{phobic\_reaction\_type([Options])} &
Indicates a type of phobic reaction, where \texttt{Options} lists the admissible
reaction options. \\
\end{longtable}
